\documentclass[a4paper,12pt]{article}

\usepackage[T2A]{fontenc}
\usepackage[utf8]{inputenc}
\usepackage[russian]{babel}
\usepackage{amsmath,amssymb,amsthm,mathtools}
\usepackage{helvet}
\usepackage{PTSans}
\renewcommand{\familydefault}{\sfdefault}
\usepackage{icomma}
\usepackage[a4paper,margin=20mm,headheight=25pt]{geometry}
\usepackage{microtype}
\usepackage{tikz}
\usetikzlibrary{calc}
\usetikzlibrary{arrows.meta,positioning,shapes.geometric}
\usepackage{tkz-euclide}
\usepackage{pgfplots}
\pgfplotsset{compat=1.18}
\usepackage{tabularx,booktabs,longtable}
\usepackage{enumitem}
\usepackage{hyperref}
\usepackage{parskip}
\usepackage{fancyhdr}
\usepackage{setspace}
\usepackage{xcolor}
\usepackage[most]{tcolorbox}

\definecolor{accent}{HTML}{008080}
\definecolor{darkaccent}{HTML}{004D4D}
\definecolor{textgray}{HTML}{333333}
\definecolor{softaccent}{HTML}{EAF4F4}
\definecolor{softgray}{HTML}{F3F5F5}

\hypersetup{hidelinks}
\color{textgray}
\onehalfspacing
\setlength{\parindent}{0pt}
\setlength{\parskip}{0.55em}
\setlist[itemize]{leftmargin=1.45em,itemsep=0.25em,topsep=0.3em}
\setlist[enumerate]{leftmargin=1.65em,itemsep=0.45em,topsep=0.35em}

\pagestyle{fancy}
\fancyhf{}
\fancyhead[L]{\small\color{darkaccent} Кирилл Зиновьев}
\fancyhead[R]{\small\color{darkaccent} Математика, 7 класс}
\fancyfoot[C]{\small\thepage}
\renewcommand{\headrulewidth}{0.3pt}
\renewcommand{\headrule}{%
  \hbox to\headwidth{%
    \color{accent}\leaders\hrule height \headrulewidth\hfill
  }%
}

\newtcolorbox{keybox}[1][]{
  enhanced,
  breakable,
  colback=softaccent,
  colframe=accent,
  boxrule=0.45pt,
  arc=1.5mm,
  left=3mm,
  right=3mm,
  top=2mm,
  bottom=2mm,
  before skip=0.7em,
  after skip=0.7em,
  #1
}

\newtcolorbox{examplebox}[1][]{
  enhanced,
  breakable,
  colback=white,
  colframe=accent!65!black,
  boxrule=0.4pt,
  arc=1.2mm,
  left=3mm,
  right=3mm,
  top=2mm,
  bottom=2mm,
  before skip=0.7em,
  after skip=0.7em,
  #1
}

\newtcolorbox{warningbox}[1][]{
  enhanced,
  breakable,
  colback=softgray,
  colframe=textgray!55,
  boxrule=0.4pt,
  arc=1.2mm,
  left=3mm,
  right=3mm,
  top=2mm,
  bottom=2mm,
  before skip=0.7em,
  after skip=0.7em,
  #1
}

\tikzset{
  flowstart/.style={
    ellipse,
    draw=accent,
    thick,
    align=center,
    minimum width=34mm,
    minimum height=10mm,
    fill=softaccent
  },
  flowaction/.style={
    rectangle,
    rounded corners=2mm,
    draw=accent,
    thick,
    align=center,
    text width=85mm,
    minimum height=11mm,
    fill=white
  },
  flowdecision/.style={
    diamond,
    aspect=2.2,
    draw=accent,
    thick,
    align=center,
    text width=47mm,
    inner sep=1.6mm,
    fill=softaccent
  },
  flowarrow/.style={
    -{Stealth[length=3mm]},
    thick,
    draw=darkaccent
  }
}

\newcommand{\sectionline}{%
  \par
  \noindent
  \color{accent}\rule{\linewidth}{0.45pt}
  \color{textgray}
  \par
}

\begin{document}

% -------------------------------------------------
% Титульный лист
% -------------------------------------------------
\begin{titlepage}
\thispagestyle{empty}
\centering

\vspace*{2.1cm}

{\Large\color{darkaccent}\textbf{Чек-лист}}\par

\vspace{0.55cm}

{\huge\bfseries
Отношения величин и задачи на прямоугольник
}\par

\vspace{0.25cm}

{\Large\color{accent}
Математическая модель текстовой задачи
}\par

\vspace{1.5cm}

{\large 7 класс}\par

\vspace{0.4cm}

{\large Кирилл Зиновьев}\par

\vfill

{\large \today}\par

\vspace{0.55cm}

{\normalsize
Лёвин Артём Александрович эксклюзивно для Кирилла Зиновьева
}\par

\vspace{1.15cm}

\begin{tikzpicture}[remember picture,overlay]
  \draw[
    accent!55,
    line width=1.1pt
  ]
    ($(current page.south west)+(18mm,10mm)$)
    .. controls
       ($(current page.south west)+(63mm,18mm)$)
       and
       ($(current page.south east)+(-63mm,5mm)$)
    ..
    ($(current page.south east)+(-18mm,10mm)$);
\end{tikzpicture}

\end{titlepage}

% -------------------------------------------------
% Основной чек-лист
% -------------------------------------------------

\section*{1. Что означает отношение величин}
\sectionline

Если стороны прямоугольника относятся как
\[
a:b=m:n,
\]
то числа $m$ и $n$ показывают количество одинаковых долей.

Для длин сторон можно записать
\[
a=mx,
\qquad
b=nx,
\qquad
x>0.
\]

Здесь $x$ --- длина одной общей доли.

\begin{keybox}
\textbf{Главная идея.}

Запись
\[
a:b=3:6
\]
не означает
\[
a=3,
\qquad
b=6.
\]

Она означает, что обе длины получаются умножением чисел $3$ и $6$ на один и тот же положительный множитель.
\end{keybox}

Например, отношению $3:6$ подходят пары длин
\[
3\text{ и }6,
\qquad
6\text{ и }12,
\qquad
9\text{ и }18,
\]
поскольку во всех случаях отношение сокращается до $1:2$.

\subsection*{Можно ли сократить отношение?}

Да. Отношения
\[
3:6
\quad\text{и}\quad
1:2
\]
равносильны.

Поэтому возможны две записи:
\[
a=3x,
\qquad
b=6x
\]
или
\[
a=y,
\qquad
b=2y.
\]

Переменные $x$ и $y$ тогда имеют разные численные значения, однако итоговые длины сторон совпадают.

\section*{2. Формулы прямоугольника}
\sectionline

Для прямоугольника со сторонами $a$ и $b$:
\[
P=2(a+b),
\qquad
S=ab.
\]

Если
\[
a=mx,
\qquad
b=nx,
\]
то
\[
P=2(mx+nx)=2(m+n)x.
\]

Площадь:
\[
S=(mx)(nx)=mnx^2.
\]

\begin{keybox}
При умножении $x$ на $x$ получаем
\[
x\cdot x=x^2.
\]

Поэтому, например,
\[
(3x)(6x)=18x^2.
\]
\end{keybox}

\section*{3. Типовой способ решения задачи}
\sectionline

Если известно отношение сторон и периметр, удобно действовать в таком порядке:

\begin{enumerate}
  \item Перевести отношение на язык переменной:
  \[
  a=mx,
  \qquad
  b=nx.
  \]

  \item Записать формулу периметра.

  \item Подставить выражения для сторон и получить уравнение относительно $x$.

  \item Найти $x$.

  \item Восстановить длины сторон.

  \item Вычислить требуемую величину, например площадь.

  \item Проверить отношение сторон и исходное условие.
\end{enumerate}

\begin{examplebox}
\textbf{Пример 1.}

Стороны прямоугольника относятся как $3:6$, периметр равен $54$ см. Найти площадь.

Пусть
\[
a=3x,
\qquad
b=6x.
\]

Тогда
\[
P=3x+6x+3x+6x=18x.
\]

По условию
\[
P=54,
\]
поэтому
\[
18x=54,
\qquad
x=3.
\]

Следовательно,
\[
a=9\text{ см},
\qquad
b=18\text{ см}.
\]

Площадь:
\[
S=9\cdot18=162\text{ см}^2.
\]
\end{examplebox}

\begin{examplebox}
\textbf{Пример 2.}

Стороны относятся как $2:3$, периметр равен $50$ см.

Записываем
\[
a=2x,
\qquad
b=3x.
\]

Тогда
\[
2x+3x+2x+3x=50,
\]
откуда
\[
10x=50,
\qquad
x=5.
\]

Стороны равны
\[
10\text{ см}
\quad\text{и}\quad
15\text{ см}.
\]

Поэтому
\[
S=150\text{ см}^2.
\]
\end{examplebox}

\begin{examplebox}
\textbf{Пример 3.}

Стороны относятся как $3:5$, периметр равен $16$ см.

Пусть
\[
a=3x,
\qquad
b=5x.
\]

Тогда
\[
3x+5x+3x+5x=16,
\]
то есть
\[
16x=16,
\qquad
x=1.
\]

Следовательно,
\[
S=(3x)(5x)=15x^2=15\text{ см}^2.
\]
\end{examplebox}

\section*{4. Математическая модель текстовой задачи}
\sectionline

Математическая модель --- это запись условий задачи с помощью переменных, формул, равенств и уравнений.

Рассмотрим движение из города $A$ в город $B$.

Часть пути пройдена по обычной дороге со скоростью $40$ км/ч, оставшаяся часть --- по шоссе со скоростью $60$ км/ч. Общий путь равен $260$ км.

Если времена движения равны $t_1$ и $t_2$, то пройденные расстояния:
\[
s_1=40t_1,
\qquad
s_2=60t_2.
\]

Поскольку весь путь равен $260$ км,
\[
40t_1+60t_2=260.
\]

\begin{warningbox}
Одного такого уравнения недостаточно для однозначного нахождения двух неизвестных $t_1$ и $t_2$.

Нужна дополнительная независимая информация, например общее время движения:
\[
t_1+t_2=T.
\]

Тогда получается система из двух уравнений.
\end{warningbox}

\subsection*{Проверка здравого смысла}

Если весь путь $260$ км проходится со скоростями от $40$ до $60$ км/ч, общее время должно лежать между
\[
\frac{260}{60}
=
\frac{13}{3}
\approx
4{,}33\text{ ч}
\]
и
\[
\frac{260}{40}
=
6{,}5\text{ ч}.
\]

Поэтому, например, общее время $4$ часа противоречило бы физическому смыслу условия.

\section*{5. Типичные ошибки}
\sectionline

\begin{itemize}
  \item Подменять отношение
  \[
  a:b=3:6
  \]
  равенствами
  \[
  a=3,
  \qquad
  b=6.
  \]

  \item Использовать разные множители для частей одного отношения. В записи
  \[
  a=3x,
  \qquad
  b=6x
  \]
  переменная $x$ одна и та же.

  \item Забывать, что
  \[
  x\cdot x=x^2.
  \]

  \item Пытаться найти площадь до определения масштаба $x$, если данных для этого ещё недостаточно.

  \item Использовать число из условия, не связывая его с формулой. Например, периметр должен войти в уравнение через
  \[
  P=2(a+b).
  \]

  \item Терять единицы измерения: длина измеряется в см, м, км; площадь --- в см$^2$, м$^2$ и т.~д.

  \item В текстовой задаче с несколькими неизвестными считать модель завершённой, когда уравнений недостаточно для однозначного решения.
\end{itemize}

% -------------------------------------------------
% Блок-схема
% -------------------------------------------------

\clearpage
\thispagestyle{plain}

\begin{center}
{\Large\bfseries\color{darkaccent}
Алгоритм: отношение сторон и известный периметр
}
\end{center}

\vspace{5mm}

\begin{center}
\begin{tikzpicture}[node distance=8mm and 16mm]

  \node[flowstart] (start) {
    Начало
  };

  \node[
    flowaction,
    below=of start
  ] (ratio) {
    Из отношения $a:b=m:n$ записать\\
    $a=mx$, $b=nx$, $x>0$
  };

  \node[
    flowaction,
    below=of ratio
  ] (perim) {
    Подставить в формулу периметра\\
    $P=2(a+b)$
  };

  \node[
    flowaction,
    below=of perim
  ] (eq) {
    Составить и решить уравнение\\
    $2(m+n)x=P$
  };

  \node[
    flowaction,
    below=of eq
  ] (sides) {
    Найти стороны\\
    $a=mx$, $b=nx$
  };

  \node[
    flowaction,
    below=of sides
  ] (target) {
    Вычислить требуемую величину\\
    например, $S=ab$
  };

  \node[
    flowdecision,
    below=10mm of target
  ] (check) {
    Условие и отношение выполняются?
  };

  \node[
    flowstart,
    below=10mm of check
  ] (finish) {
    Ответ
  };

  \node[
    flowaction,
    right=24mm of check,
    text width=55mm
  ] (revise) {
    Проверить арифметику,\\
    уравнение и подстановку
  };

  \draw[flowarrow]
    (start) -- (ratio);

  \draw[flowarrow]
    (ratio) -- (perim);

  \draw[flowarrow]
    (perim) -- (eq);

  \draw[flowarrow]
    (eq) -- (sides);

  \draw[flowarrow]
    (sides) -- (target);

  \draw[flowarrow]
    (target) -- (check);

  \draw[flowarrow]
    (check)
    --
    node[left,font=\small]{да}
    (finish);

  \draw[flowarrow]
    (check.east)
    --
    node[above,font=\small]{нет}
    (revise.west);

  \draw[flowarrow]
    (revise.north)
    |-
    ([xshift=8mm]eq.east)
    --
    (eq.east);

\end{tikzpicture}
\end{center}

\clearpage

% -------------------------------------------------
% Тренировочный блок
% -------------------------------------------------

\section*{Тренировочный блок --- Путь А}
\sectionline

\textbf{Правило работы.}

Оформляйте решение последовательно: введите общую долю $x$, составьте уравнение, найдите $x$, затем вычислите требуемую величину.

\begin{enumerate}
  \item
  Стороны прямоугольника относятся как $2:3$, а его периметр равен $40$ см. Найдите площадь прямоугольника.

  \item
  Стороны прямоугольника относятся как $3:5$, а его периметр равен $64$ см. Найдите площадь прямоугольника.

  \item
  Стороны прямоугольника относятся как $4:7$, а его периметр равен $66$ см. Найдите длины сторон и площадь прямоугольника.

  \item
  Стороны прямоугольника относятся как $5:8$, а его периметр равен $104$ см. Найдите площадь прямоугольника.

  \item
  Автомобиль проехал из города $A$ в город $B$ всего $260$ км.

  Часть пути он ехал по обычной дороге со скоростью $40$ км/ч, оставшуюся часть --- по шоссе со скоростью $60$ км/ч.

  Общее время движения составило $5$ часов.

  Сколько времени автомобиль ехал по обычной дороге и сколько по шоссе?
\end{enumerate}

\vfill

\begin{keybox}
\textbf{Самопроверка перед завершением работы:}

все числа из условия использованы по смыслу;

уравнение соответствует условию;

найденное значение $x$ положительно;

отношение сторон сохраняется;

единицы измерения указаны.
\end{keybox}

% -------------------------------------------------
% Ответы
% -------------------------------------------------

\clearpage

\section*{Ответы}
\sectionline

\renewcommand{\arraystretch}{1.35}

\begin{table}[ht]
\centering

\begin{tabularx}{\linewidth}{@{}>{\bfseries}cX@{}}
\toprule
№ & Ответ \\
\midrule
1 & $96\text{ см}^2$ \\

2 & $240\text{ см}^2$ \\

3 & стороны $12$ см и $21$ см; площадь $252\text{ см}^2$ \\

4 & $640\text{ см}^2$ \\

5 & по обычной дороге --- $2$ ч; по шоссе --- $3$ ч \\
\bottomrule
\end{tabularx}

\end{table}

\vfill

\begin{center}
\small\color{darkaccent}
Лёвин Артём Александрович эксклюзивно для Кирилла Зиновьева
\end{center}

\end{document}